%!TeX root=csp-proof.tex
\section{The algorithm}\label{sec:algorithm}

The algorithm is \cite{ZhukJACM}'s. It is stated here because target
\textup{(T1)} of \S\ref{sec:target} is a proposition \emph{about} it and
cannot be written down otherwise.

\subsection{\textsc{Solve}}

\begin{verbatim}
function Solve(I):
    repeat
        I := ForceConsistency(I)
        if I = false: return "No"
        if some D_x has a strong subset B:
            I := ReduceDomain(I, x, B)
    until nothing changed
    return SolveLinear(I)
\end{verbatim}

\textsc{ForceConsistency} enforces cycle-consistency and irreducibility,
returning \texttt{false} if some domain empties. The reduction step is licensed
by Theorem~\ref{thm:reductions-safe}: restricting to a strong subset cannot
destroy solvability. When no strong subset remains, Lemma~\ref{lem:ubiquity}
together with Lemma~\ref{lem:linear-on-top} gives, for each domain of size at
least $2$, a congruence $\sigma_{x}$ with $D_{x}/\sigma_{x}\cong\mathbb Z_{q_{1}}
\times\dots\times\mathbb Z_{q_{n_{x}}}$; choose $\sigma_{x}$ minimal. That case
is \textsc{SolveLinear}.

\subsection{\textsc{SolveLinear}}

Write $\phi\colon\mathbb Z_{p_{1}}\times\dots\times\mathbb Z_{p_{k}}\to
D_{x_{1}}/\sigma_{x_{1}}\times\dots\times D_{x_{m}}/\sigma_{x_{m}}$ for a linear
map and
\[
  \phi^{-1}(\mathcal I)=\{\bar a:\mathcal I\text{ has a solution in }\phi(\bar a)\}.
\]
Computing $\phi^{-1}(\mathcal I)$ would solve $\mathcal I$; the algorithm cannot
do that, but it can do four things:

\begin{description}\tightlist
\item[\textup{(p0)}] decide $\bar a\in\phi^{-1}(\mathcal I)$, by reducing each
  domain to the corresponding block and recursing on a smaller domain;
\item[\textup{(p1)}] decide $\phi^{-1}(\mathcal I)=\mathbb Z_{p_{1}}\times\dots
  \times\mathbb Z_{p_{k}}$, by testing $\bar 0$ first and then the $k$ unit
  vectors with \textup{(p0)}. The set $\phi^{-1}(\mathcal I)$ is in general an
  affine \emph{coset}, not a subgroup; testing $\bar 0$ first is what makes it
  one, after which membership of the coordinate generators implies fullness;
\item[\textup{(p2)}] compute $\phi^{-1}(\mathcal I)$ when $\mathcal I$ is not
  linked, by splitting into linked pieces, recursing, and taking the union;
\item[\textup{(p3)}] compute $\phi^{-1}(\mathcal I)$ when it is empty or of
  codimension one, by finding $\bar a\notin\phi^{-1}(\mathcal I)$, then for each
  $i$ a $b_{i}$ with $\bar a[i\mapsto b_{i}]\in\phi^{-1}(\mathcal I)$ if one
  exists, and reading off the equation
  $\sum_{i\in J}(y_{i}-a_{i})/(b_{i}-a_{i})=1$ over the $i$ for which $b_{i}$
  exists. The displayed quotient is meaningful only within one prime: the
  codimension-one equation has a single target prime $p$, and only coordinates
  over $\mathbb F_{p}$ can carry a nonzero coefficient, so the computation is
  organised prime by prime (Caveat~\ref{haz:step2-linear-algebra}).
\end{description}

\begin{verbatim}
function SolveLinear(I):
    m := dim(D_{x_1}/s_1 x ... x D_{x_n}/s_n)
    phi := a bijective linear map Z_{p_1}x...xZ_{p_m} -> that product
    while phi^-1(I) != Z_{p_1}x...xZ_{p_m}:              // (p1)
        I' := I
        for C in I':                            // drop what can be dropped
            if phi^-1(I' \ {C}) != Z_{p_1}x...xZ_{p_m}:   // (p1)
                I' := I' \ {C}
        F := phi^-1(I')                                // (p2) or (p3)
        if F = {}: return "No"
        m := dim F ;  psi := a bijective linear map onto F ;  phi := phi o psi
    return "Yes"
\end{verbatim}

The inner loop makes $\mathcal I'$ minimal with $\phi^{-1}(\mathcal I')$ not
full. At that point Theorem~\ref{thm:codim-one} applies --- its hypothesis (vii)
is exactly minimality --- so $\phi^{-1}(\mathcal I')$ is empty or of codimension
one, and \textup{(p3)} computes it; or $\mathcal I'$ is not linked and
\textup{(p2)} does. Since $\phi^{-1}(\mathcal I)\subseteq\phi^{-1}(\mathcal I')$,
the image of $\phi$ shrinks while still containing all solutions. The dimension
strictly decreases, so the loop terminates.

\subsection{What \textup{(T1)} has to say}

\begin{hazard}[Two obstacles between the pseudocode and a function]
\label{haz:algorithm-obstacles}
\begin{enumerate}\alphlabels\tightlist
\item \emph{Termination is not structural.} \textsc{Solve} recurses through
  \textup{(p0)} on strictly smaller domains, and \textsc{SolveLinear} loops on
  a strictly decreasing dimension. Neither is a structural recursion; both need
  an explicit well-founded measure, and the two interleave ---
  \textsc{SolveLinear} calls \textsc{Solve} on a smaller domain, which calls
  \textsc{SolveLinear} again. The measure is lexicographic:
  $\bigl(\sum_{x}|D_{x}|,\ \dim\phi\bigr)$.
\item \emph{Several steps quantify over objects that are only finitely many by
  accident.} ``If some $D_{x}$ has a strong subset $B$'' quantifies over subsets
  of $D_{x}$, over congruences on $D_{x}$, and --- through absorption --- over
  \emph{term operations}, of which there are infinitely many. For the types the
  algorithm actually searches for, the search is finite for a reason that needs
  no general arity bound: $\TBA$ has a binary witness by definition, and a
  central subuniverse has a ternary one by Lemma~\ref{lem:central-ternary}.
\end{enumerate}
\end{hazard}

\subsection{\textsc{Solve} as a function}

The pseudocode is not a function: it makes unconstrained choices, and its
recursion is not structural. Both are fixed here, so that
Theorem~\ref{thm:t1} has a subject.

\begin{definition}[The call set and its measure]\label{def:solve-measure}
A \emph{call} is either $\mathsf S(\mathcal I)$, for $\mathcal I$ an instance
of $\CSP(\Gamma)$ with domains $D_{x}$, or $\mathsf L(\mathcal I,\phi)$, where
in addition $\phi$ is an injective linear map
$\mathbb Z_{p_{1}}\times\dots\times\mathbb Z_{p_{k}}\to\prod_{x}
D_{x}/\sigma_{x}$ with the $\sigma_{x}$ the minimal congruences of
\S\ref{sec:algorithm}. Put $N(\mathcal I)=\sum_{x}|D_{x}|$ and
\[
  \mu\bigl(\mathsf S(\mathcal I)\bigr)=
    \bigl(N(\mathcal I),\,N(\mathcal I)+1\bigr),
  \qquad
  \mu\bigl(\mathsf L(\mathcal I,\phi)\bigr)=
    \bigl(N(\mathcal I),\,\dim\phi\bigr),
\]
ordered lexicographically on $\mathbb N^{2}$. Since $\dim\phi\le N(\mathcal I)$
always, a call $\mathsf S(\mathcal I)$ is strictly above every
$\mathsf L(\mathcal I,\phi)$ with the same instance.
\end{definition}

\begin{definition}[$\textsc{Solve}_{\mathrm{alg}}$]\label{def:solve-alg}
Fix once and for all a linear order on the finite set of pairs $(x,B)$ with
$B\subseteq D_{x}$, and one on the finite set of candidate maps $\phi$; ``the
least'' below means least in these orders, which is what turns each choice in
the pseudocode into a value. Define $\textsc{Solve}_{\mathrm{alg}}$ on calls by
well-founded recursion on $\mu$:

\begin{description}\tightlist
\item[$\mathsf S(\mathcal I)$.] Put $\mathcal I'=
  \textsc{ForceConsistency}(\mathcal I)$. If some domain of $\mathcal I'$ is
  empty, return \texttt{false}. If some $D'_{x}$ has a proper nonempty subset
  $B$ with $B\subt{T}D'_{x}$ for some $T\in\{\TBA,\TC,\TD\}$, take the least
  such $(x,B)$ and return
  $\textsc{Solve}_{\mathrm{alg}}\bigl(\mathsf S(\textsc{ReduceDomain}(\mathcal
  I',x,B))\bigr)$. Otherwise return $\textsc{Solve}_{\mathrm{alg}}
  \bigl(\mathsf L(\mathcal I',\phi_{0})\bigr)$ for the least admissible
  $\phi_{0}$.
\item[$\mathsf L(\mathcal I,\phi)$.] Run the body of \textsc{SolveLinear},
  reading each test $\bar a\in\phi^{-1}(\mathcal I)$ of \textup{(p0)} as the
  value of $\textsc{Solve}_{\mathrm{alg}}$ at $\mathsf S$ of the instance
  $\mathcal I$ with each domain cut to the block $\phi(\bar a)_{x}$, and each
  iteration of the loop as the value at $\mathsf L(\mathcal I,\phi\circ\psi)$.
\end{description}
\end{definition}

\begin{lemma}[The recursion is well founded]\label{lem:solve-terminates}
Every recursive call made by Definition~\ref{def:solve-alg} is at a strictly
smaller $\mu$, so $\textsc{Solve}_{\mathrm{alg}}$ is a total function into
$\{\texttt{true},\texttt{false}\}$.
\end{lemma}

\begin{proof}
Four kinds of call. \textsc{ForceConsistency} only shrinks domains, so it does
not increase $N$; the reduction step replaces some $D'_{x}$ by a proper subset,
strictly decreasing $N$. The call from $\mathsf S(\mathcal I')$ to
$\mathsf L(\mathcal I',\phi_{0})$ leaves $N$ fixed and drops the second
component, by the inequality in Definition~\ref{def:solve-measure}. The
\textup{(p0)} calls cut every domain to a block of $\sigma_{x}$; each
$\sigma_{x}$ is proper at some variable, since $\mathsf L$ is reached only when
some domain has size at least $2$, so $N$ strictly decreases. The loop calls
leave $N$ fixed and strictly decrease $\dim\phi$, as
\S\ref{sec:algorithm} records.
\end{proof}

\begin{theorem}[Target \textup{(T1)}]\label{thm:t1}
For every instance $\mathcal I$ of $\CSP(\Gamma)$,
\[
  \textsc{Solve}_{\mathrm{alg}}\bigl(\mathsf S(\mathcal I)\bigr)=\texttt{true}
  \iff \mathcal I\ \text{has a solution}.
\]
\end{theorem}

\begin{hazard}[What Definition~\ref{def:solve-alg} still owes]\label{haz:t1-names-solve}
Theorem~\ref{thm:t1} is about \emph{this} function, and that is the whole of
its content: the weaker reading --- ``there exists a Boolean-valued function
deciding satisfiability'' --- follows for a fixed finite template by exhaustive
search and says nothing. Three things are needed before the definition is
airtight, and none is deep.

\begin{enumerate}\alphlabels\tightlist
\item \textsc{ForceConsistency} and \textsc{ReduceDomain} are used as given
  total functions. The first must be exhibited as a fixpoint iteration that
  terminates because it only shrinks domains; the second is immediate.
\item The search in $\mathsf S$ ranges over subsets $B$ and over the witnesses
  making $B\subt{T}D'_{x}$. It is finite for the reason in
  Caveat~\ref{haz:algorithm-obstacles}(b) --- $\TBA$ has a binary witness by
  definition and $\TC$ a ternary one by Lemma~\ref{lem:central-ternary}, and
  $\TD$ quantifies over congruences --- but the bound has to be stated for the
  ``least such $(x,B)$'' to be well defined.
\item Only $\TBA$, $\TC$ and $\TD$ are searched for, not $\TS$; that is
  legitimate because $\subt{\TS}$ is witnessed by a subset carrying the other
  two types, but the correctness proof must use
  Theorem~\ref{thm:reductions-safe} in the form that covers exactly the types
  searched.
\end{enumerate}

Theorem~\ref{thm:t1} itself is stated and not proved here; see
\S\ref{sec:status}.
\end{hazard}

Given that, Theorem~\ref{thm:t1} mentions no machine, is provable from
\S\ref{sec:main-statements}, and is the honest content of ``the algorithm
works'': that one particular, highly non-obvious algorithm --- the one whose
existence is the dichotomy --- is correct.

\begin{theorem}[Target \textup{(T2)}; stated, not proved]\label{thm:t2}
For each finite $\Gamma$ there are a program computing
$\textsc{Solve}_{\mathrm{alg}}$ and a
polynomial $p_{\Gamma}$ bounding its running time on instances of
$\CSP(\Gamma)$ in the size of the instance.
\end{theorem}

\begin{hazard}[\textup{(T2)} is not a wrapper]\label{haz:t2-scope}
Theorem~\ref{thm:t2} is the only statement of the tractable half that needs a
cost model, and it is not a corollary of Theorem~\ref{thm:t1}: it requires
$\textsc{Solve}_{\mathrm{alg}}$ to be exhibited as a program with a proved step
bound. Note
also that the statement is non-uniform. Because $\Gamma$ is fixed, the searches
over subsets of $A$, over congruences on $A$, and over term operations of
bounded arity are constant-time --- but ``constant'' means compiled away into a
program whose \emph{size} depends on $|A|$. So $\textsc{Solve}_{\mathrm{alg}}$
is a family with a polynomial for each $\Gamma$, and a bound
uniform in $\Gamma$ would be a different statement, and a false one.
\end{hazard}
